
\section*{CHƯƠNG 5 - KẾT LUẬN VÀ HƯỚNG PHÁT TRIỂN}
\addcontentsline{toc}{section}{\numberline{}CHƯƠNG 5 - KẾT LUẬN VÀ HƯỚNG PHÁT TRIỂN}
\setcounter{section}{5}
\setcounter{subsection}{0}
\setcounter{figure}{0}
\setcounter{table}{0}

\subsection{Tổng kết kết quả đạt được}

\subsubsection{Về mặt lý thuyết}
Đề tài đã hoàn thành các mục tiêu nghiên cứu lý thuyết:

\begin{itemize}
    \item \textbf{Sentiment Analysis:} Đã nghiên cứu và nắm vững khái niệm phân tích cảm xúc, các phương pháp tiếp cận (Lexicon-based, Machine Learning, Deep Learning), và ứng dụng thực tế trong phân tích ý kiến khách hàng.
    
    \item \textbf{Natural Language Processing (NLP):} Đã hiểu rõ các kỹ thuật xử lý ngôn ngữ tự nhiên cơ bản:
    \begin{itemize}
        \item Text preprocessing: Tokenization, lowercasing, removing special characters
        \item Text normalization: Stemming, lemmatization
        \item Feature extraction: Bag of Words, TF-IDF
        \item Sentiment scoring: Polarity và Subjectivity
    \end{itemize}
    
    \item \textbf{YouTube Data API v3:} Đã nắm vững cách sử dụng YouTube API để thu thập dữ liệu:
    \begin{itemize}
        \item Cơ chế authentication và API keys
        \item Quota management và rate limiting
        \item Pagination với nextPageToken
        \item Xử lý lỗi và failover mechanism
    \end{itemize}
    
    \item \textbf{TextBlob Library:} Đã hiểu rõ cách hoạt động của TextBlob:
    \begin{itemize}
        \item Pattern-based sentiment analysis
        \item Polarity score (-1 đến 1)
        \item Subjectivity score (0 đến 1)
        \item Ưu nhược điểm so với các phương pháp khác
    \end{itemize}
    
    \item \textbf{Data Mining Pipeline:} Đã nắm vững quy trình khai phá dữ liệu 6 bước:
    \begin{itemize}
        \item Data Collection: Thu thập dữ liệu từ nguồn
        \item Data Cleaning: Làm sạch và chuẩn hóa
        \item Data Labeling: Gán nhãn tự động
        \item Experimentation: Thực nghiệm và phân tích
        \item Application: Xây dựng ứng dụng
        \item Deep Analysis: Phân tích sâu kết quả
    \end{itemize}
\end{itemize}

\subsubsection{Về mặt kỹ thuật}
Đề tài đã triển khai thành công hệ thống hoàn chỉnh:

\begin{itemize}
    \item \textbf{Thu thập dữ liệu:} Đã xây dựng hệ thống thu thập tự động với:
    \begin{itemize}
        \item 4 API keys dự phòng với cơ chế failover
        \item Xử lý pagination để thu thập lượng lớn dữ liệu
        \item Error handling và retry mechanism
        \item Thu thập thành công 100,000 bình luận
    \end{itemize}
    
    \item \textbf{Chuẩn hóa dữ liệu:} Đã triển khai quy trình làm sạch 8 bước chi tiết:
    \begin{itemize}
        \item Loại bỏ URLs, mentions, hashtags, emails
        \item Loại bỏ ký tự đặc biệt và số điện thoại
        \item Chuẩn hóa khoảng trắng
        \item Chuyển về chữ thường
        \item Lọc bình luận rỗng, quá ngắn, trùng lặp
        \item Thống kê chi tiết từng bước xử lý
    \end{itemize}
    
    \item \textbf{Gán nhãn tự động:} Đã triển khai hệ thống phân loại 3 mức:
    \begin{itemize}
        \item Positive/Negative/Neutral với ngưỡng linh hoạt
        \item 3 mức độ tin cậy: High/Medium/Low
        \item 5 cột output: score, subjectivity, label, confidence, description
        \item Xử lý 10,000 mẫu với thời gian hợp lý (8 phút)
    \end{itemize}
    
    \item \textbf{Trực quan hóa:} Đã tạo 6 loại biểu đồ:
    \begin{itemize}
        \item Bar Chart: Phân bố số lượng theo nhãn
        \item Pie Chart: Tỷ lệ phần trăm các nhãn
        \item Scatter Plot: Độ dài vs Cảm xúc, Likes vs Cảm xúc
        \item Box Plot: Phân bố điểm theo nhãn
        \item Horizontal Bar: Độ dài trung bình theo nhãn
        \item Correlation Heatmap: Ma trận tương quan
    \end{itemize}
    
    \item \textbf{Ứng dụng demo:} Đã xây dựng ứng dụng phân tích văn bản mới với 2 chế độ:
    \begin{itemize}
        \item \textbf{Chế độ 1:} Test với 5 mẫu văn bản có sẵn (demo tự động)
        \item \textbf{Chế độ 2:} Chế độ tương tác - cho phép người dùng tự nhập văn bản không giới hạn
        \item Input: Văn bản tiếng Anh bất kỳ
        \item Output: Nhãn, điểm số, độ tin cậy, thống kê chi tiết
        \item Hiển thị với emoji và màu sắc
        \item Chế độ interactive cho người dùng tự nhập và phân tích
        \item Validation và error handling đầy đủ
    \end{itemize}
\end{itemize}

\subsubsection{Về mặt kết quả}
Đề tài đã đạt được các kết quả quan trọng:

\begin{itemize}
    \item \textbf{Dữ liệu thu thập:}
    \begin{itemize}
        \item 100,000 bình luận từ video Avengers: Endgame
        \item Sau làm sạch: 92,521 bình luận hợp lệ (92.5\%)
        \item Phân tích chi tiết: 10,000 mẫu ngẫu nhiên
    \end{itemize}
    
    \item \textbf{Phân bố cảm xúc:}
    \begin{itemize}
        \item Positive: 62.3\% - Phần lớn khán giả đánh giá tích cực
        \item Negative: 18.8\% - Một số ý kiến tiêu cực
        \item Neutral: 18.9\% - Ý kiến trung lập
    \end{itemize}
    
    \item \textbf{Độ chính xác:}
    \begin{itemize}
        \item Tổng thể: 83\% (kiểm tra thủ công 100 mẫu)
        \item Positive: 88.3\% - Phân loại tốt nhất
        \item Negative: 85.0\% - Phân loại tốt
        \item Neutral: 70.0\% - Phân loại khó khăn hơn
    \end{itemize}
    
    \item \textbf{Hiệu năng:}
    \begin{itemize}
        \item Thời gian tổng: \textasciitilde 60 phút cho toàn bộ pipeline
        \item Thu thập: 45 phút (100K bình luận)
        \item Chuẩn hóa: 3 phút (92K bình luận)
        \item Gán nhãn: 8 phút (10K mẫu)
        \item Trực quan hóa: 2 phút
    \end{itemize}
    
    \item \textbf{Phát hiện insights:}
    \begin{itemize}
        \item Bình luận tích cực có xu hướng dài hơn (132 ký tự vs 115-118)
        \item Bình luận tích cực nhận nhiều likes hơn (correlation 0.234)
        \item Độ dài văn bản không ảnh hưởng nhiều đến cảm xúc (correlation 0.123)
        \item Phim Avengers: Endgame được đón nhận rất tích cực (62.3\% positive)
    \end{itemize}
\end{itemize}

\subsection{Đánh giá ưu nhược điểm}

\subsubsection{Ưu điểm của giải pháp}

\textbf{Về phương pháp:}
\begin{itemize}
    \item \textbf{Unsupervised Learning:} Không cần dữ liệu training được gán nhãn thủ công, tiết kiệm thời gian và chi phí
    \item \textbf{Lexicon-based:} TextBlob sử dụng từ điển cảm xúc, dễ hiểu và giải thích kết quả
    \item \textbf{Nhanh chóng:} Xử lý 10,000 mẫu chỉ trong 8 phút
    \item \textbf{Đơn giản:} Dễ triển khai, không cần GPU hay tài nguyên lớn
\end{itemize}

\textbf{Về hệ thống:}
\begin{itemize}
    \item \textbf{Tự động hóa:} Toàn bộ pipeline chạy tự động từ đầu đến cuối
    \item \textbf{Failover mechanism:} 4 API keys dự phòng đảm bảo thu thập liên tục
    \item \textbf{Chi tiết:} 8 bước làm sạch với thống kê từng bước
    \item \textbf{Trực quan:} 6 loại biểu đồ giúp hiểu rõ dữ liệu
    \item \textbf{Mở rộng:} Dễ dàng thay đổi video, tăng số mẫu, hoặc thêm tính năng
\end{itemize}

\textbf{Về kết quả:}
\begin{itemize}
    \item \textbf{Độ chính xác tốt:} 83\% cho phương pháp unsupervised
    \item \textbf{Insights hữu ích:} Phát hiện xu hướng và mối quan hệ giữa các đặc trưng
    \item \textbf{Ứng dụng thực tế:} Demo có thể phân tích văn bản mới ngay lập tức
\end{itemize}

\subsubsection{Nhược điểm và hạn chế}

\textbf{Về phương pháp:}
\begin{itemize}
    \item \textbf{Lexicon-based limitations:}
    \begin{itemize}
        \item Không hiểu ngữ cảnh phức tạp
        \item Không nhận diện sarcasm (mỉa mai)
        \item Không xử lý tốt slang và từ viết tắt
        \item Không học từ dữ liệu (không adaptive)
    \end{itemize}
    
    \item \textbf{Chỉ hỗ trợ tiếng Anh:} TextBlob hoạt động tốt nhất với tiếng Anh, hạn chế với ngôn ngữ khác
    
    \item \textbf{Neutral khó phân loại:} Độ chính xác thấp (70\%) do ranh giới mờ giữa Neutral và Positive/Negative nhẹ
\end{itemize}

\textbf{Về dữ liệu:}
\begin{itemize}
    \item \textbf{Giới hạn API:} Quota 10,000 units/day/key hạn chế tốc độ thu thập
    \item \textbf{Chỉ phân tích 1 video:} Kết quả chỉ đại diện cho Avengers: Endgame
    \item \textbf{Bias:} Dữ liệu YouTube có thể bị bias (người xem tích cực thường comment nhiều hơn)
    \item \textbf{Spam và bot:} Có thể có bình luận spam hoặc bot không được lọc hết
\end{itemize}

\textbf{Về hệ thống:}
\begin{itemize}
    \item \textbf{Không có validation set:} Chỉ kiểm tra thủ công 100 mẫu, chưa có test set lớn
    \item \textbf{Không có model training:} Không sử dụng Machine Learning để cải thiện độ chính xác
    \item \textbf{Thiếu tính năng:} Chưa có phân tích aspect-based, emotion detection, topic modeling
\end{itemize}

\subsection{Bài học kinh nghiệm}

\subsubsection{Kiến thức kỹ thuật}
\begin{itemize}
    \item \textbf{API management quan trọng:} Cần có nhiều API keys dự phòng và xử lý lỗi tốt để đảm bảo thu thập dữ liệu liên tục.
    
    \item \textbf{Data cleaning là then chốt:} Chất lượng dữ liệu sau làm sạch ảnh hưởng trực tiếp đến kết quả phân tích. Cần làm sạch kỹ lưỡng và có thống kê chi tiết.
    
    \item \textbf{TextBlob có giới hạn:} Phù hợp cho prototype và ứng dụng đơn giản, nhưng cần phương pháp nâng cao hơn (ML/DL) cho độ chính xác cao.
    
    \item \textbf{Visualization giúp hiểu dữ liệu:} Biểu đồ giúp phát hiện patterns và insights mà số liệu thô không thể hiện.
    
    \item \textbf{Neutral là thách thức:} Phân loại Neutral khó hơn Positive/Negative do ranh giới mờ và ngữ cảnh phức tạp.
\end{itemize}

\subsubsection{Kỹ năng thực hành}
\begin{itemize}
    \item \textbf{Modular code:} Chia code thành các function riêng biệt giúp dễ debug và maintain.
    
    \item \textbf{Progress tracking:} Hiển thị tiến trình xử lý giúp người dùng biết hệ thống đang làm gì.
    
    \item \textbf{Error handling:} Xử lý lỗi tốt giúp hệ thống robust và không bị crash giữa chừng.
    
    \item \textbf{Documentation:} Comment code và tạo README giúp người khác hiểu và sử dụng dễ dàng.
    
    \item \textbf{Testing:} Test với nhiều trường hợp khác nhau (positive, negative, neutral, edge cases) để đảm bảo hệ thống hoạt động đúng.
\end{itemize}

\subsubsection{Quản lý dự án}
\begin{itemize}
    \item \textbf{Chia nhỏ công việc:} Chia dự án thành 6 bước rõ ràng giúp dễ quản lý và theo dõi tiến độ.
    
    \item \textbf{Làm việc nhóm:} Phân công rõ ràng giữa 2 thành viên, thường xuyên sync và review code.
    
    \item \textbf{Version control:} Sử dụng Git để quản lý code và dễ dàng rollback khi cần.
    
    \item \textbf{Deadline management:} Đặt deadline cho từng bước và buffer time cho các vấn đề không lường trước.
\end{itemize}

\subsection{Hướng phát triển trong tương lai}

\subsubsection{Cải thiện độ chính xác}

\textbf{Machine Learning Models:}
\begin{itemize}
    \item Triển khai các mô hình ML: Naive Bayes, SVM, Random Forest
    \item So sánh hiệu năng với TextBlob
    \item Sử dụng labeled dataset (IMDB, Twitter Sentiment) để training
    \item Cross-validation và hyperparameter tuning
\end{itemize}

\textbf{Deep Learning Models:}
\begin{itemize}
    \item Triển khai LSTM, GRU cho sequence modeling
    \item Sử dụng pre-trained models: BERT, RoBERTa, DistilBERT
    \item Fine-tuning trên domain-specific data (movie reviews)
    \item Đánh giá trade-off giữa độ chính xác và tốc độ
\end{itemize}

\textbf{Ensemble Methods:}
\begin{itemize}
    \item Kết hợp nhiều models (TextBlob + ML + DL)
    \item Voting hoặc stacking để cải thiện độ chính xác
    \item Weighted ensemble dựa trên confidence score
\end{itemize}

\subsubsection{Mở rộng tính năng}

\textbf{Aspect-Based Sentiment Analysis:}
\begin{itemize}
    \item Phân tích cảm xúc theo khía cạnh: plot, acting, visual effects, music
    \item Sử dụng dependency parsing và named entity recognition
    \item Tạo báo cáo chi tiết về từng aspect
\end{itemize}

\textbf{Emotion Detection:}
\begin{itemize}
    \item Mở rộng từ 3 nhãn (Pos/Neg/Neu) sang 6-8 emotions: joy, sadness, anger, fear, surprise, disgust
    \item Sử dụng emotion lexicons (NRC, EmoLex)
    \item Trực quan hóa emotion distribution
\end{itemize}

\textbf{Topic Modeling:}
\begin{itemize}
    \item Áp dụng LDA (Latent Dirichlet Allocation) để phát hiện topics
    \item Phân tích cảm xúc theo từng topic
    \item Tìm hiểu khán giả quan tâm đến điều gì nhất
\end{itemize}

\textbf{Temporal Analysis:}
\begin{itemize}
    \item Phân tích cảm xúc theo thời gian (ngày, tuần, tháng)
    \item Phát hiện xu hướng thay đổi cảm xúc
    \item So sánh cảm xúc trước và sau sự kiện (premiere, awards)
\end{itemize}

\subsubsection{Mở rộng dữ liệu}

\textbf{Multi-Video Analysis:}
\begin{itemize}
    \item Phân tích nhiều videos của cùng một phim
    \item So sánh cảm xúc giữa trailer, review, reaction videos
    \item Tổng hợp insights từ nhiều nguồn
\end{itemize}

\textbf{Multi-Platform:}
\begin{itemize}
    \item Thu thập dữ liệu từ Twitter, Reddit, IMDB
    \item So sánh cảm xúc giữa các platforms
    \item Phân tích cross-platform sentiment
\end{itemize}

\textbf{Multi-Language:}
\begin{itemize}
    \item Hỗ trợ tiếng Việt, tiếng Trung, tiếng Nhật
    \item Sử dụng multilingual models (mBERT, XLM-R)
    \item So sánh cảm xúc giữa các quốc gia
\end{itemize}

\subsubsection{Cải thiện hệ thống}

\textbf{Web Application:}
\begin{itemize}
    \item Xây dựng web app với Flask/Django
    \item Giao diện user-friendly cho người dùng không kỹ thuật
    \item Upload video URL và nhận kết quả phân tích
    \item Dashboard với biểu đồ interactive (Plotly, D3.js)
\end{itemize}

\textbf{Real-time Analysis:}
\begin{itemize}
    \item Stream comments real-time từ YouTube
    \item Phân tích và cập nhật kết quả liên tục
    \item Alert khi phát hiện sentiment shift đột ngột
\end{itemize}

\textbf{API Service:}
\begin{itemize}
    \item Xây dựng RESTful API cho sentiment analysis
    \item Cho phép các ứng dụng khác tích hợp
    \item Rate limiting và authentication
    \item Documentation với Swagger/OpenAPI
\end{itemize}

\textbf{Database Integration:}
\begin{itemize}
    \item Lưu trữ dữ liệu trong database (PostgreSQL, MongoDB)
    \item Query và filter dữ liệu hiệu quả
    \item Historical data analysis
\end{itemize}

\subsubsection{Ứng dụng thực tế}

\textbf{Business Intelligence:}
\begin{itemize}
    \item Phân tích phản hồi khách hàng về sản phẩm/dịch vụ
    \item Theo dõi brand reputation trên social media
    \item Phát hiện crisis sớm và phản ứng kịp thời
\end{itemize}

\textbf{Movie Industry:}
\begin{itemize}
    \item Dự đoán box office dựa trên sentiment
    \item Phân tích phản hồi về trailer để điều chỉnh marketing
    \item So sánh sentiment giữa các phim cùng thể loại
\end{itemize}

\textbf{Academic Research:}
\begin{itemize}
    \item Nghiên cứu về cultural differences trong cảm xúc
    \item Phân tích ảnh hưởng của social media đến opinion
    \item Publish papers về sentiment analysis methods
\end{itemize}

\subsection{Kết luận}

Đề tài "Phân lớp Sentiment Analysis với YouTube API" đã hoàn thành các mục tiêu đề ra:

\textbf{Về lý thuyết:} Đã nghiên cứu và nắm vững các khái niệm về Sentiment Analysis, NLP, YouTube API, TextBlob, và quy trình Data Mining 6 bước.

\textbf{Về kỹ thuật:} Đã thiết kế và triển khai thành công hệ thống hoàn chỉnh với 6 bước: Thu thập dữ liệu (100K bình luận) → Chuẩn hóa (8 bước làm sạch) → Gán nhãn (10K mẫu) → Thực nghiệm (6 biểu đồ) → Ứng dụng (demo) → Phân tích sâu (insights).

\textbf{Về kết quả:} Đã đạt được độ chính xác 83\% với phương pháp unsupervised, phát hiện được insights hữu ích (62.3\% positive, bình luận tích cực dài hơn và nhận nhiều likes hơn), và xây dựng được ứng dụng demo hoạt động tốt.

\textbf{Về ứng dụng:} Hệ thống có thể áp dụng thực tế để phân tích phản hồi khách hàng, theo dõi brand reputation, dự đoán xu hướng, và hỗ trợ ra quyết định kinh doanh.

Đề tài cũng mở ra nhiều hướng phát triển trong tương lai: Cải thiện độ chính xác với ML/DL, mở rộng tính năng (aspect-based, emotion detection, topic modeling), mở rộng dữ liệu (multi-video, multi-platform, multi-language), và xây dựng web application với real-time analysis.

Qua quá trình thực hiện đề tài, nhóm đã tích lũy được nhiều kiến thức và kỹ năng quý báu về Data Mining, NLP, Sentiment Analysis, Python programming, API integration, data visualization, và project management. Những kiến thức này sẽ là nền tảng vững chắc cho việc học tập và nghiên cứu sâu hơn trong lĩnh vực Data Science và AI.

Nhóm xin chân thành cảm ơn TS. Trần Thanh Phước đã hướng dẫn và hỗ trợ nhóm trong suốt quá trình thực hiện đề tài. Nhóm cũng xin cảm ơn Khoa Công Nghệ Thông Tin, Trường Đại học Tôn Đức Thắng đã tạo điều kiện để nhóm hoàn thành đề tài này.

\vspace{1cm}

\textbf{Thành viên nhóm:}
\begin{itemize}
    \item Đỗ Tài Khải Nguyên - MSSV: 52300230
    \item Tống Phương Nam - MSSV: 52300130
\end{itemize}

\textbf{Giảng viên hướng dẫn:}
\begin{itemize}
    \item TS. Trần Thanh Phước
\end{itemize}

\textbf{Lớp:} DH52CS

\textbf{Năm học:} 2023-2027
